\documentclass[11pt]{article}
\usepackage[margin=1.5cm,paperwidth=16cm,paperheight=20cm]{geometry}
\usepackage{amsmath,amssymb,amsthm}
\usepackage{microtype}
\usepackage[T1]{fontenc}
\usepackage{lmodern}
\usepackage{xcolor}
\usepackage{booktabs}
\usepackage{parskip}

\definecolor{accentblue}{RGB}{30,80,160}

\pagestyle{empty}

\begin{document}

\begin{center}
  {\large\bfseries\color{accentblue} RetroDiff: Training Objective Derivation}
\end{center}

\medskip
\noindent
Diffusion models are trained by maximising a variational lower bound on the log-likelihood.
Applying Jensen's inequality to the intractable marginal $\log p_\theta(x)$ gives the
\emph{evidence lower bound} (ELBO):

\begin{align}
  \log p_\theta(x)
    &\;\geq\;
    \mathbb{E}_{q(x_{1:T}|x)}\!\left[
      \log\frac{p_\theta(x_{0:T})}{q(x_{1:T}|x)}
    \right]
    \;=\; -\mathcal{L}_{\mathrm{ELBO}}
  \label{eq:elbo}
\end{align}

\noindent
where $q$ is the fixed forward noising process,
$p_\theta$ is the learned reverse process parameterised by $\theta$,
and $x_{1:T}$ denotes the latent variables across all $T$ diffusion steps.

\medskip
\noindent
Expanding the joint ratio in \eqref{eq:elbo} via the Markov structure of both
$p_\theta$ and $q$ yields a telescoping decomposition into a reconstruction term
and a sum of per-step KL divergences:

\begin{align}
  \mathcal{L}_{\mathrm{ELBO}}
    &= \underbrace{\mathbb{E}_q[-\log p_\theta(x_0|x_1)]}_{\text{reconstruction}}
    \;+\;
    \sum_{t=2}^{T}
    \underbrace{
      \mathrm{KL}\!\left(
        q(x_{t-1}|x_t, x) \,\|\, p_\theta(x_{t-1}|x_t)
      \right)
    }_{\text{step-}t\text{ denoising}}
  \label{eq:decomp}
\end{align}

\noindent
Each KL term in \eqref{eq:decomp} can be computed in closed form because the
forward posterior $q(x_{t-1}|x_t, x)$ is Gaussian; minimising these terms trains
the network to reverse each noising step.

\medskip
\noindent
\textbf{RetroDiff's contribution.}
Standard diffusion models generate $x$ from noise alone.
RetroDiff conditions every denoising step on a set
$R(x)=\{r_1,\ldots,r_k\}$ of $k$ exemplars retrieved from the training corpus
based on $x$.
Substituting $\mathcal{L}_{\mathrm{ELBO}}(x\,|\,R(x))$
for $\mathcal{L}_{\mathrm{ELBO}}(x)$ in the per-sample objective and taking an
expectation over the data distribution defines the full training loss:

\begin{align}
  \mathcal{L}_{\mathrm{RetroDiff}}(\theta)
    &= \mathbb{E}_{x \sim p_{\mathrm{data}}}\!\left[
        \mathcal{L}_{\mathrm{ELBO}}\!\left(x \,\big|\, R(x)\right)
      \right]
  \label{eq:retro}
\end{align}

\noindent
Equation~\eqref{eq:retro} is the central objective: the model must learn to
denoise $x$ given both the noisy intermediate $x_t$ and the retrieved exemplar
context, making retrieval a first-class component of the generative process
rather than a post-hoc conditioning signal.

\bigskip
\begin{center}
\small
\renewcommand{\arraystretch}{1.3}
\begin{tabular}{cl}
  \toprule
  \textbf{Symbol} & \textbf{Meaning} \\
  \midrule
  $x$ & clean data sample \\
  $x_t$ & noised version of $x$ at diffusion step $t$ \\
  $R(x)$ & set of $k$ retrieved exemplars conditioned on $x$ \\
  $\theta$ & model parameters \\
  $q$ & forward noising distribution \\
  $p_\theta$ & learned reverse distribution \\
  $T$ & total diffusion steps \\
  \bottomrule
\end{tabular}
\end{center}

\end{document}
